% GENERATED FILE. Edit canonical lessons, then run npm run generate:books.
% canonical-source-hash: fnv1a64:387981f5933e8975
% canonical-lessons: TE-C434-vyavastha, TE-C434-vidhanam, TE-C434-sthiti, TE-C434-sandarbham, TE-C434-erpadu

\chapter{Systems and Situations}
\label{ch:te-systems-and-situations}
\hlchaptermodality{434}

\begin{chapteropening}
\textbf{By the end of this chapter:} \emph{I can talk about systems, policies, states of affairs and occasions.}

Complete the last lesson of chapter 434: I can talk about systems, policies, states of affairs and occasions.
\end{chapteropening}

\section[\texorpdfstring{\te{వ్యవస్థ} (vyavastha)}{వ్యవస్థ (vyavastha)}]{\te{వ్యవస్థ} (vyavastha) — a system; an institution}

\label{lesson:TE-C434-vyavastha}



\begin{quote}
Before the new one: say the Telugu for dominance, then the Telugu for to be deprived of.
\end{quote}

\subsection*{You'll want to know: \te{వ్యవస్థ}}
\textbf{\te{వ్యవస్థ}} — \emph{vyavastha} — \textquotedblleft{}a system; an institution\textquotedblright{}.

\begin{quote}
\te{మన} \te{విద్యా} \te{వ్యవస్థలో} \te{మార్పులు} \te{రావాలి}.
\end{quote}

\begin{quote}
\emph{mana vidyā vyavasthalō mārpulu rāvāli.}
\end{quote}

\begin{quote}
Changes must come in our education system.
\end{quote}

Plural: \te{వ్యవస్థలు}, \emph{vyavasthalu}.

Plural \te{వ్యవస్థలు}.

\begin{cousinweb}
\textbf{\te{వ్యవస్థ}} (\emph{vyavastha}) is borrowed from Sanskrit \emph{vyavasthā}, \textquotedblleft{}respective difference; abiding in one place, steadiness\textquotedblright{} (Monier-Williams); Telugu uses it for a settled order, a system.
\end{cousinweb}

\begin{grammarlens}[title={What you've built}]
One more word for saying what happens if.
\end{grammarlens}

\subsection*{Your turn}
\begin{itemize}
\raggedright
  \item \emph{Say it:} \emph{vyavastha}
  \item \emph{Say it:} \emph{vyavastha}, once more
  \item \emph{Say it:} \emph{mana vidyā vyavasthalō mārpulu rāvāli.} — Changes must come in our education system.
  \item \emph{Recall:} say the Telugu for dominance, then the Telugu for to be deprived of, then say \emph{vyavastha} again
\end{itemize}

\begin{tcolorbox}[breakable,colback=teal!4,colframe=teal!35!black,title={Before you move on}]
What does \te{వ్యవస్థ} mean? (\textquotedblleft{}A system; an institution\textquotedblright{}.) Say it once more.
\end{tcolorbox}
\section[\texorpdfstring{\te{విధానం} (vidhānaṁ)}{విధానం (vidhānaṁ)}]{\te{విధానం} (vidhānaṁ) — a policy; a method}

\label{lesson:TE-C434-vidhanam}



\begin{quote}
Before the new one: say the Telugu for to be deprived of, then the Telugu for a system.
\end{quote}

\subsection*{You'll want to know: \te{విధానం}}
\textbf{\te{విధానం}} — \emph{vidhānaṁ} — \textquotedblleft{}a policy; a method\textquotedblright{}.

\begin{quote}
\te{ప్రభుత్వం} \te{కొత్త} \te{విద్యా} \te{విధానం} \te{ప్రకటించింది}.
\end{quote}

\begin{quote}
\emph{prabhutvaṁ kotta vidyā vidhānaṁ prakaṭiñcindi.}
\end{quote}

\begin{quote}
The government announced a new education policy.
\end{quote}

Plural: \te{విధానాలు}, \emph{vidhānālu}.

Also a way of doing: \te{పని} \te{చేసే} \te{విధానం}, the way of working.

\begin{cousinweb}
\textbf{\te{విధానం}} (\emph{vidhānaṁ}) is borrowed from Sanskrit \emph{vidhāna}, \textquotedblleft{}arrangement, order\textquotedblright{} (CDIAL 11754).
\end{cousinweb}

\begin{grammarlens}[title={What you've built}]
One more word for saying what happens if.
\end{grammarlens}

\subsection*{Your turn}
\begin{itemize}
\raggedright
  \item \emph{Say it:} \emph{vidhānaṁ}
  \item \emph{Say it:} \emph{vidhānaṁ}, once more
  \item \emph{Say it:} \emph{prabhutvaṁ kotta vidyā vidhānaṁ prakaṭiñcindi.} — The government announced a new education policy.
  \item \emph{Recall:} say the Telugu for to be deprived of, then the Telugu for a system, then say \emph{vidhānaṁ} again
\end{itemize}

\begin{tcolorbox}[breakable,colback=teal!4,colframe=teal!35!black,title={Before you move on}]
What does \te{విధానం} mean? (\textquotedblleft{}A policy; a method\textquotedblright{}.) Say it once more.
\end{tcolorbox}
\section[\texorpdfstring{\te{స్థితి} (sthiti)}{స్థితి (sthiti)}]{\te{స్థితి} (sthiti) — a state of affairs; a state}

\label{lesson:TE-C434-sthiti}



\begin{quote}
Before the new one: say the Telugu for a system, then the Telugu for a policy.
\end{quote}

\subsection*{You'll want to know: \te{స్థితి}}
\textbf{\te{స్థితి}} — \emph{sthiti} — \textquotedblleft{}a state of affairs; a state\textquotedblright{}.

\begin{quote}
\te{రోగి} \te{స్థితి} \te{ఇప్పుడు} \te{బాగుంది}.
\end{quote}

\begin{quote}
\emph{rōgi sthiti ippuḍu bāgundi.}
\end{quote}

\begin{quote}
The patient is in a good state now.
\end{quote}

\te{పరిస్థితి}, a situation, which you know, is the same word with \te{పరి}-.

\begin{cousinweb}
\textbf{\te{స్థితి}} (\emph{sthiti}) is borrowed from Sanskrit \emph{sthiti}, \textquotedblleft{}standing upright, continued existence\textquotedblright{} (CDIAL 13769), from \emph{sthā}, \textquotedblleft{}to stand\textquotedblright{}.
\end{cousinweb}

\begin{grammarlens}[title={What you've built}]
One more word for saying what happens if.
\end{grammarlens}

\subsection*{Your turn}
\begin{itemize}
\raggedright
  \item \emph{Say it:} \emph{sthiti}
  \item \emph{Say it:} \emph{sthiti}, once more
  \item \emph{Say it:} \emph{rōgi sthiti ippuḍu bāgundi.} — The patient is in a good state now.
  \item \emph{Recall:} say the Telugu for a system, then the Telugu for a policy, then say \emph{sthiti} again
\end{itemize}

\begin{tcolorbox}[breakable,colback=teal!4,colframe=teal!35!black,title={Before you move on}]
What does \te{స్థితి} mean? (\textquotedblleft{}A state of affairs; a state\textquotedblright{}.) Say it once more.
\end{tcolorbox}
\section[\texorpdfstring{\te{సందర్భం} (sandarbhaṁ)}{సందర్భం (sandarbhaṁ)}]{\te{సందర్భం} (sandarbhaṁ) — an occasion; a context}

\label{lesson:TE-C434-sandarbham}



\begin{quote}
Before the new one: say the Telugu for a policy, then the Telugu for a state of affairs.
\end{quote}

\subsection*{You'll want to know: \te{సందర్భం}}
\textbf{\te{సందర్భం}} — \emph{sandarbhaṁ} — \textquotedblleft{}an occasion; a context\textquotedblright{}.

\begin{quote}
\te{ఆ} \te{మాట} \te{ఏ} \te{సందర్భంలో} \te{అన్నారు}?
\end{quote}

\begin{quote}
\emph{ā māṭa ē sandarbhaṁlō annāru?}
\end{quote}

\begin{quote}
In what context was that said?
\end{quote}

Plural: \te{సందర్భాలు}, \emph{sandarbhālu}.

\te{సందర్భంగా}, on the occasion of: \te{పుట్టినరోజు} \te{సందర్భంగా}, on the occasion of a birthday.

\begin{cousinweb}
\textbf{\te{సందర్భం}} (\emph{sandarbhaṁ}) is borrowed from Sanskrit \emph{saṁdarbha}, \textquotedblleft{}stringing or binding together; a literary composition\textquotedblright{} (Monier-Williams); Telugu uses it for the setting of a word or an event.
\end{cousinweb}

\begin{grammarlens}[title={What you've built}]
One more word for saying what happens if.
\end{grammarlens}

\subsection*{Your turn}
\begin{itemize}
\raggedright
  \item \emph{Say it:} \emph{sandarbhaṁ}
  \item \emph{Say it:} \emph{sandarbhaṁ}, once more
  \item \emph{Say it:} \emph{ā māṭa ē sandarbhaṁlō annāru?} — In what context was that said?
  \item \emph{Recall:} say the Telugu for a policy, then the Telugu for a state of affairs, then say \emph{sandarbhaṁ} again
\end{itemize}

\begin{tcolorbox}[breakable,colback=teal!4,colframe=teal!35!black,title={Before you move on}]
What does \te{సందర్భం} mean? (\textquotedblleft{}An occasion; a context\textquotedblright{}.) Say it once more.
\end{tcolorbox}
\section[\texorpdfstring{\te{ఏర్పడు} (ērpaḍu)}{ఏర్పడు (ērpaḍu)}]{\te{ఏర్పడు} (ērpaḍu) — to be formed; to take shape}

\label{lesson:TE-C434-erpadu}



\begin{quote}
Before the new one: say the Telugu for a state of affairs, then the Telugu for an occasion.
\end{quote}

\subsection*{You'll want to know: \te{ఏర్పడు}}
\textbf{\te{ఏర్పడు}} — \emph{ērpaḍu} — \textquotedblleft{}to be formed; to take shape\textquotedblright{}.

\begin{quote}
\te{ఆ} \te{సంఘం} \te{పది} \te{సంవత్సరాల} \te{కిందట} \te{ఏర్పడింది}.
\end{quote}

\begin{quote}
\emph{ā saṅghaṁ padi saṁvatsarāla kindaṭa ērpaḍindi.}
\end{quote}

\begin{quote}
That association was formed ten years ago.
\end{quote}

Past (it …): \te{ఏర్పడింది}, \emph{ērpaḍindi}.

Groups, habits and ideas take shape: \te{అలవాటు} \te{ఏర్పడింది}, a habit formed. Its past is shown in the third person, \te{ఏర్పడింది}; \te{ఏర్పడ్డాను}, I was formed, is not said.

\begin{cousinweb}
\textbf{\te{ఏర్పడు}} (\emph{ērpaḍu}) is native Dravidian: DEDR 905 lists Telugu \emph{ērpaḍu} beside \emph{ēlu}, \textquotedblleft{}to accept, to admit\textquotedblright{}, with -\te{పడు}, the verb \emph{paḍu} (DEDR 3853).
\end{cousinweb}

\begin{grammarlens}[title={What you've built}]
One more word for saying what happens if.
\end{grammarlens}

\subsection*{Your turn}
\begin{itemize}
\raggedright
  \item \emph{Say it:} \emph{ērpaḍu}
  \item \emph{Say it:} \emph{ērpaḍu}, once more
  \item \emph{Say it:} \emph{ā saṅghaṁ padi saṁvatsarāla kindaṭa ērpaḍindi.} — That association was formed ten years ago.
  \item \emph{Recall:} say the Telugu for a state of affairs, then the Telugu for an occasion, then say \emph{ērpaḍu} again
\end{itemize}

\begin{tcolorbox}[breakable,colback=teal!4,colframe=teal!35!black,title={Before you move on}]
What does \te{ఏర్పడు} mean? (\textquotedblleft{}To be formed; to take shape\textquotedblright{}.) Say it once more.
\end{tcolorbox}
